% !TeX root = ../strategy_report.tex
\chapter{Dual Pipeline Architecture}
\label{ch:dual_pipeline}

A common anti-pattern in machine learning deployments is the monolithic pipeline, where neural network training sequences and data indexing tasks are tightly coupled into a single execution graph. The TaxonVision-MLOps architecture explicitly rejects this model. Instead, as visualised in the comprehensive architectural overview following Chapter 1, it establishes a dual-pipeline topology that decouples the weight-centric classification engine from the data-centric Multimodal Knowledge Graph (MMKG). These pipelines maintain independent lifecycles, converging exclusively at the inference serving boundary.

\section{Asymmetric Iteration and Execution Decoupling}
\label{sec:decoupling}

The necessity for architectural decoupling stems directly from the fundamentally asymmetric nature of biological modelling operations.

\textbf{Pipeline 1: The Classification Engine.} The primary neural backbone (for example, an optimised BioCLIP-2 \cite{stevens2024bioclip} or DINOv3 transformer) is computationally expensive to train but relatively static in production. Refining its weights using Class-Balanced Loss \cite{cui2019classbalanced} requires intensive GPU Tensor Core utilisation, distributed gradient synchronisation, and static ONNX Post-Training Quantisation (PTQ) calibration. However, once an optimal checkpoint is compiled, the model architecture may remain unchanged for weeks or months.

\textbf{Pipeline 2: The Taxonomic MMKG.} Conversely, the underlying citizen-science data stream is highly volatile. The pipeline must ingest continuous drops of new iNaturalist research-grade observations. More critically, biological taxonomies are frequently subjected to scientific realignments -- species are split into novel taxonomic branches or lumped together under unified nomenclature. Pipeline 2 is an I/O-bound vector-indexing workflow responsible for parsing DarwinCore Parquet manifests, extracting zero-shot organism crops, filtering open-source licences, and appending dense embeddings to the Hierarchical Navigable Small World (HNSW) graph.

Forcing Pipeline 1 to re-execute a heavy neural network training sweep merely because Pipeline 2 updated an exemplar image for a specific beetle species wastes immense computational resources. Furthermore, physically separating these domains enforces strict failure isolation. If a malformed DarwinCore archive from an upstream provider induces a schema violation, Pipeline 2 will halt securely. However, Pipeline 1 remains entirely unimpeded, preserving our ability to compile, test, and deploy critical hot-fixes to the classification logic.

\section{Automated CI/CD Orchestration \& Registry Promotion}
\label{sec:ci_cd_binding}

\begin{figure}[htbp]
\centering
\begin{tikzpicture}[
    node distance=1.5cm and 2cm,
    box/.style={draw, rectangle, rounded corners, minimum width=2.5cm, minimum height=1cm, align=center, font=\small},
    arrow/.style={-{Stealth[scale=1.2]}, thick},
    dashed_arrow/.style={-{Stealth[scale=1.2]}, dashed, thick},
    cloud/.style={draw, ellipse, minimum height=1cm, align=center, font=\small, fill=gray!10},
    db/.style={draw, cylinder, shape border rotate=90, aspect=0.25, minimum height=1.2cm, minimum width=1.5cm, align=center, font=\small}
]

% Nodes
\node[box] (git) {Git Push\\(Codebase)};
\node[box, right=of git] (gh_actions) {GitHub\\(Gates \& Webhook)};
\node[box, right=of gh_actions] (runner) {Jenkins Agent Pod\\(k3s, RTX 5070 Ti)};
\node[db, below=of runner] (dvc) {DVC S3 Remote\\(DagsHub)};
\node[box, below=of gh_actions] (training) {Training \& ONNX\\Quantisation};
\node[box, left=of training] (mlflow) {MLflow Registry\\(Safety Gate)};
\node[cloud, below=of mlflow] (k8s) {Kubernetes Ingress\\(Serving \& RAG)};
\node[box, right=of k8s] (prom) {Prometheus\\(Drift Alert)};

% Edges
\draw[arrow] (git) -- node[above, font=\footnotesize] {Trigger} (gh_actions);
\draw[arrow] (gh_actions) -- node[above, font=\footnotesize] {Webhook} (runner);
\draw[arrow] (runner) -- node[right, font=\footnotesize] {Pull Data} (dvc);
\draw[arrow] (runner) -- (training);
\draw[arrow] (training) -- node[above, font=\footnotesize, align=left] {Log \&\\Register} (mlflow);
\draw[arrow] (mlflow) -- node[left, font=\footnotesize] {Deploy} (k8s);
\draw[arrow] (k8s) -- node[above, font=\footnotesize] {Telemetry} (prom);
\draw[dashed_arrow] (prom) to[out=30, in=230] node[left, font=\footnotesize, align=left] {Retrain\\Trigger} (runner);

\end{tikzpicture}
\caption{Dual-Pipeline CI/CD Architecture Flow mapping code and data states to production Kubernetes deployments.}
\label{fig:cicd_flow}
\end{figure}

In order to transition from manual experimentation to a Level 3 MLOps deployment, the independent pipelines must be strictly governed by Continuous Integration and Continuous Deployment (CI/CD) triggers.

Two cooperating layers implement these triggers. When engineers push codebase modifications to the central repository, GitHub Actions executes lightweight static analysis and logical unit tests on CPU-only runner pods. In parallel, a signed GitHub webhook reaches a self-hosted Jenkins controller through an outbound-only Cloudflare Tunnel. The controller owns no executors of its own: it schedules each build as an ephemeral agent pod inside the k3s cluster of the bare-metal host and removes the pod on completion. Credentials for the tracking infrastructure, in particular the DagsHub access token, are held in the Jenkins credential store and are exposed to the build as environment variables for the duration of the run only.

As the training pipeline initialises, it is specified to bind the MLflow telemetry to the exact Git commit that triggered the run, thereby ensuring cryptographic provenance between the source code and the resulting model weights:
\begin{minted}{python}
mlflow.set_tag("git.commit", os.environ.get("GIT_COMMIT"))
\end{minted}

Crucially, the CI/CD pipeline enforces an automated epistemic safety gate prior to deployment. Before any compiled model is allowed to enter the registry, a standalone invariant test loads the candidate ONNX graph and evaluates it against a holdout calibration dataset. The script formally asserts that the empirical marginal coverage of the Split Conformal Prediction set, $|C|$, meets or exceeds the mathematical bound $1 - \alpha$ (for example, $95\%$ coverage, where $\alpha=0.05$). If the empirical coverage falls below this strict threshold, the pipeline throws a fatal assertion error, blocking the Git merge and effectively preventing the uncalibrated model from reaching production.

If the safety gate is cleared, the system programmatically interacts with the MLflow API. The candidate is compared with the model carrying the \texttt{Production} alias by validation Macro PR-AUC: a better candidate is marked as \texttt{Challenger}, whereas an inferior one is archived. The promotion of a \texttt{Challenger} to \texttt{Production} remains a deliberate step after deployment testing, so that at most two competing models are active at any time:
\begin{minted}{python}
client.set_registered_model_alias(
    name="taxon_vision_classifier",
    alias="Challenger",
    version=mv.version,
)
\end{minted}

\section{Automated Training on the Main Branch}
\label{sec:main_training}

On the main branch, the Jenkins pipeline continues beyond the quality gates and executes the training stage on the bare-metal host. The stage is deliberately placed after the documentation build, so that no GPU time is spent on a revision that would fail one of the earlier gates. It proceeds as follows:
\begin{enumerate}
    \item The frozen GPU environment is installed from the dependency lock file, and the versioned dataset is pulled from the DagsHub DVC remote.
    \item The classification head is trained on the embeddings of the configured backbone (DINOv3 by default) within an epoch budget of 15 and a mini-batch size of 64, under the process-level VRAM cap of Section~\ref{sec:gpu_partitioning}. A Bayesian hyperparameter search (Section~\ref{sec:hyperparameter_opt}) can be requested beforehand by a command-line switch.
    \item The ONNX graph is re-exported, and the resulting head checkpoint and graph are pushed to the DVC remote.
    \item A rolling restart of the serving deployment loads the fresh artefacts without interrupting service: the Kubernetes rolling-update strategy admits one additional pod and no unavailable pod, and a readiness probe withholds traffic until the new pod has loaded the model.
\end{enumerate}

\section{Convergence at the Serving Boundary}
\label{sec:serving_convergence}

The two distinct pipelines exist in absolute isolation until the precise moment of user inference. The FastAPI serving application acts as the convergence node, simultaneously loading the INT8 quantised ONNX graph from Pipeline 1 into memory, whilst mounting the live HNSW vector index compiled by Pipeline 2.

When a field photograph is submitted to the API, it is passed through the fast-path ONNX backbone. If the epistemic safety layer determines the classification is decisive -- meaning the conformal prediction set yields exactly one candidate ($|C| = 1$) -- the pipeline returns the prediction immediately, bypassing the vector database entirely. If, however, the model exhibits severe ambiguity ($|C| > 1$), the serving application lazy-loads Pipeline 2's HNSW index, executes a single-stage filtered traversal, and retrieves the top-$k$ reference exemplar crops in order to synthesise a robust, comparative visual context.

\section{Low-Latency Interpretability via Forward-Hooked CAM}
\label{sec:cam_shortcut}

Providing diagnostic transparency to human reviewers requires visual interpretability tools such as Class Activation Mapping (CAM). While modern extensions like Grad-CAM are popular, they require a computationally expensive backward pass during inference. Executing gradient backpropagation on a production server destroys the sub-$25\text{ ms}$ latency Service Level Agreement (SLA) and prevents the model from being deployed in highly optimized, forward-only runtimes like INT8 ONNX.

To achieve visual interpretability without violating strict operational latency constraints, the architecture bypasses the backward pass entirely by leveraging the properties of the linear classification head.

During the initial ONNX forward pass, the system utilizes execution hooks to capture the spatial feature maps $A^k$ emitted by the final convolutional or attention layer before global average pooling is applied. Because the final classification head is a simple linear transformation, the importance weight $w_k^c$ of feature map $k$ for a target class $c$ is exactly equivalent to the learned weight matrix of that linear layer.

The class activation map $L^c_{\text{CAM}}$ is synthesized through a simple linear combination of the forward spatial feature maps and the static classification weights:
$$L^c_{\text{CAM}} = \text{ReLU} \left( \sum_k w_k^c A^k \right)$$
Where:
\begin{itemize}
    \item $L^c_{\text{CAM}}$ is the resulting two-dimensional activation heat map for class $c$, which is subsequently upsampled to match the resolution of the original image.
    \item $A^k$ denotes the $k$-th channel of the spatial feature map tensor extracted during the standard forward pass.
    \item $w_k^c$ is the scalar weight connecting feature channel $k$ to the output logit for class $c$, retrieved instantly from the static model weights.
    \item $\text{ReLU}$ (Rectified Linear Unit) is applied to discard negative spatial influences, retaining only the pixels that positively contributed to the classification of class $c$.
\end{itemize}

By relying solely on cached forward activations and static linear weights, this shortcut generates high-fidelity spatial heat maps essentially for free, completing the interpretability generation in under $2\text{ ms}$ while fully preserving the forward-only, quantization-friendly execution paradigm.
